% Options for packages loaded elsewhere
\PassOptionsToPackage{unicode}{hyperref}
\PassOptionsToPackage{hyphens}{url}
\documentclass[
]{article}
\usepackage{xcolor}
\usepackage[margin=1in]{geometry}
\usepackage{amsmath,amssymb}
\setcounter{secnumdepth}{-\maxdimen} % remove section numbering
\usepackage{iftex}
\ifPDFTeX
  \usepackage[T1]{fontenc}
  \usepackage[utf8]{inputenc}
  \usepackage{textcomp} % provide euro and other symbols
\else % if luatex or xetex
  \usepackage{unicode-math} % this also loads fontspec
  \defaultfontfeatures{Scale=MatchLowercase}
  \defaultfontfeatures[\rmfamily]{Ligatures=TeX,Scale=1}
\fi
\usepackage{lmodern}
\ifPDFTeX\else
  % xetex/luatex font selection
\fi
% Use upquote if available, for straight quotes in verbatim environments
\IfFileExists{upquote.sty}{\usepackage{upquote}}{}
\IfFileExists{microtype.sty}{% use microtype if available
  \usepackage[]{microtype}
  \UseMicrotypeSet[protrusion]{basicmath} % disable protrusion for tt fonts
}{}
\makeatletter
\@ifundefined{KOMAClassName}{% if non-KOMA class
  \IfFileExists{parskip.sty}{%
    \usepackage{parskip}
  }{% else
    \setlength{\parindent}{0pt}
    \setlength{\parskip}{6pt plus 2pt minus 1pt}}
}{% if KOMA class
  \KOMAoptions{parskip=half}}
\makeatother
\setlength{\emergencystretch}{3em} % prevent overfull lines
\providecommand{\tightlist}{%
  \setlength{\itemsep}{0pt}\setlength{\parskip}{0pt}}
\usepackage{bookmark}
\IfFileExists{xurl.sty}{\usepackage{xurl}}{} % add URL line breaks if available
\urlstyle{same}
\hypersetup{
  hidelinks,
  pdfcreator={LaTeX via pandoc}}

\author{}
\date{}

\begin{document}

\section{La Historia Detrás de la Tesis: De Cajas Negras a Explicaciones
Confiables}\label{la-historia-detruxe1s-de-la-tesis-de-cajas-negras-a-explicaciones-confiables}

\emph{Una narrativa sencilla sobre Inteligencia Artificial Explicable
(XAI), pensada para no expertos y útil como estructura de la defensa
doctoral.}

\begin{center}\rule{0.5\linewidth}{0.5pt}\end{center}

\subsection{Acto 1: La Pregunta Inicial --- El Robot Mago y la Caja
Negra}\label{acto-1-la-pregunta-inicial-el-robot-mago-y-la-caja-negra}

Imagina que en tu escuela contratan a un robot superinteligente para
tomar decisiones importantes. Por ejemplo: - Quién entra al equipo de
fútbol. - Quién recibe una beca de estudios. - Quién necesita ayuda
extra en matemáticas.

El robot analiza miles de datos en un segundo y anuncia su veredicto:
\textbf{``Aceptado''} o \textbf{``Rechazado''}.

Cuando le preguntas al robot: \emph{«Oye, ¿por qué me rechazaste?»}, el
robot no responde nada. Se queda en silencio. Su cerebro es un laberinto
gigante de millones de números que nadie puede leer a simple vista.

En el mundo real de la tecnología, a estos modelos de Inteligencia
Artificial los llamamos \textbf{``Cajas Negras''}. Se usan en hospitales
para detectar enfermedades, en bancos para dar préstamos y en sistemas
legales.

Pero surge un problema grave: \textbf{no podemos ni debemos confiar a
ciegas en un robot mago si no sabemos cómo toma sus decisiones.} ¿Qué
pasa si el robot está cometiendo un error? ¿O si está discriminando a
alguien por una razón equivocada?

\begin{center}\rule{0.5\linewidth}{0.5pt}\end{center}

\subsection{Acto 2: El Dilema --- ¿Quién Explica al
Explicador?}\label{acto-2-el-dilema-quiuxe9n-explica-al-explicador}

Para resolver este problema, los científicos inventaron a los
\textbf{``Detectives de IA''} (en ciencia los llamamos \emph{métodos de
explicabilidad post-hoc}, como LIME, SHAP, Anchors o DiCE).

El trabajo de un detective de IA es mirar a la caja negra por fuera,
hacerle preguntas rápidas y decirle a las personas en español sencillo:
\textgreater{} \emph{«El robot te rechazó la beca por tus notas en
historia, no por tu edad».}

¡Parecía la solución perfecta! Pero pronto los científicos descubrieron
un nuevo dilema: \textbf{los detectives a veces se contradicen entre
sí.}

Si le pides una explicación al detective \textbf{LIME}, te dice una
cosa. Si le pides la explicación al detective \textbf{SHAP} sobre la
misma persona, ¡te dice algo diferente!

Aquí nace la \textbf{motivación central de mi Tesis de Doctorado}:
\textgreater{} \emph{Si las herramientas que nos explican la IA tampoco
son perfectas y se contradicen\ldots{} ¿cómo sabemos cuál dice la
verdad? ¿Quién evalúa a los evaluadores?}

\begin{center}\rule{0.5\linewidth}{0.5pt}\end{center}

\subsection{Acto 3: Las 6 Piezas del Rompecabezas (Los Artículos A al
F)}\label{acto-3-las-6-piezas-del-rompecabezas-los-artuxedculos-a-al-f}

Para resolver este rompecabezas, organizamos nuestra investigación en
seis etapas conectadas entre sí, como los capítulos de una historia de
detectives:

\begin{verbatim}
[Mago Caja Negra] ──> [Detectives XAI] ──> [¿En quién confiar?]
                                                   │
  ┌────────────────────────────────────────────────┴────────────────────────────────┐
  │                                                                                 │
  ▼                                                                                 ▼
[Papel A: El Reglamento] ──> [Papel B: El Duelo] ──> [Papel C: Resistencia] ──> [Papel D: Caso Real]
  (Protocolo FOM-7)          (LIME vs SHAP)          (Escalabilidad)           (UCI Adult)
                                                                                    │
  ┌─────────────────────────────────────────────────────────────────────────────────┘
  │
  ▼
[Papel E: La Balanza] ──> [Papel F: Juicio Humano] ──> [¡Transparencia Defendible!]
  (Trade-offs)              (Alineación con personas)
\end{verbatim}

\begin{center}\rule{0.5\linewidth}{0.5pt}\end{center}

\subsubsection{1. Papel A: El Reglamento para los Detectives (Protocolo
FOM-7)}\label{papel-a-el-reglamento-para-los-detectives-protocolo-fom-7}

\begin{itemize}
\tightlist
\item
  \textbf{Estado:} Publicado en \emph{Revista de Investigación
  Multidisciplinaria Iberoamericana (RIMI)} (2026, DOI:
  10.69850/rimi.vi3.307).
\item
  \textbf{La idea simple:} Antes de poner a trabajar a los detectives,
  necesitamos un reglamento oficial de 7 pasos para asegurarnos de que
  las pruebas sean válidas. Lo llamamos el protocolo \textbf{FOM-7}.
\item
  \textbf{Lo que descubrimos:} Creamos una forma justa de medir 5
  propiedades clave:

  \begin{enumerate}
  \def\labelenumi{\arabic{enumi}.}
  \tightlist
  \item
    \textbf{Fidelidad:} ¿La explicación dice la verdad sobre lo que hizo
    la IA?
  \item
    \textbf{Estabilidad:} Si le preguntamos dos veces lo mismo,
    ¿mantiene su respuesta o cambia de opinión?
  \item
    \textbf{Parsimonia:} ¿La explicación es corta y fácil de entender?
  \item
    \textbf{Coste:} ¿Cuánto tiempo y cálculo le toma al detective
    responder?
  \item
    \textbf{Brecha de fidelidad:} ¿Qué tanto se equivoca al simplificar
    el problema?
  \end{enumerate}
\end{itemize}

\begin{center}\rule{0.5\linewidth}{0.5pt}\end{center}

\subsubsection{2. Papel B: El Gran Duelo --- LIME vs SHAP a
Escala}\label{papel-b-el-gran-duelo-lime-vs-shap-a-escala}

\begin{itemize}
\tightlist
\item
  \textbf{Estado:} Enviado a \emph{CLEI Electronic Journal (CLEIej)}
  (2026-10-04, Subm. \#1196).
\item
  \textbf{La idea simple:} Pusimos a competir frente a frente a los dos
  detectives más famosos del mundo en cientos de situaciones diferentes.
\item
  \textbf{Lo que descubrimos:}

  \begin{itemize}
  \tightlist
  \item
    \textbf{SHAP} es como \emph{Sherlock Holmes}: es extremadamente
    cuidadoso, nunca cambia de opinión (muy estable) y da respuestas muy
    precisas, pero es muy lento y costoso.
  \item
    \textbf{LIME} es como un \emph{detective veloz}: responde rapidísimo
    y da explicaciones muy breves, pero si el viento sopla (pequeños
    cambios en los datos), a veces cambia un poco su versión.
  \end{itemize}
\end{itemize}

\begin{center}\rule{0.5\linewidth}{0.5pt}\end{center}

\subsubsection{3. Papel C: La Prueba de Resistencia y
Escalabilidad}\label{papel-c-la-prueba-de-resistencia-y-escalabilidad}

\begin{itemize}
\tightlist
\item
  \textbf{Estado:} Listo para envío a \emph{Tecnología en Marcha}
  (Edición Especial IA, borrador de 15 pág., Zenodo v0.12.0).
\item
  \textbf{La idea simple:} ¿Qué pasa cuando el problema se vuelve 10
  veces más grande o los datos tienen ruido? ¿Los detectives resisten la
  presión?
\item
  \textbf{Lo que descubrimos:} Evaluamos cómo se comportan los
  explicadores cuando aumentamos el tamaño de la muestra y la
  complejidad de los modelos. Identificamos los puntos exactos donde las
  explicaciones dejan de ser confiables si no se controlan los
  parámetros.
\end{itemize}

\begin{center}\rule{0.5\linewidth}{0.5pt}\end{center}

\subsubsection{4. Papel D: El Examen en el Mundo Real (Benchmark UCI
Adult)}\label{papel-d-el-examen-en-el-mundo-real-benchmark-uci-adult}

\begin{itemize}
\tightlist
\item
  \textbf{Estado:} Listo para envío a \emph{Tecnología en Marcha}
  (Edición Especial IA, solicitud de cuenta pendiente).
\item
  \textbf{La idea simple:} Vamos a probar a todos los detectives en un
  caso real importante: predecir los ingresos de las personas con el
  conjunto de datos \emph{UCI Adult Income}.
\item
  \textbf{Lo que descubrimos:} Evaluamos 5 familias de modelos (árboles,
  redes neuronales, regresión logística, etc.) cruzados con 4 tipos de
  explicadores. Demostramos en la práctica cómo se aplican las 7 puertas
  del protocolo FOM-7 para obtener evidencia auditable en lugar de
  suposiciones.
\end{itemize}

\begin{center}\rule{0.5\linewidth}{0.5pt}\end{center}

\subsubsection{5. Papel E: La Balanza de las Decisiones (Trade-offs y
Gobernanza)}\label{papel-e-la-balanza-de-las-decisiones-trade-offs-y-gobernanza}

\begin{itemize}
\tightlist
\item
  \textbf{Estado:} Enviado a \emph{Computación y Sistemas (CIC-IPN)}
  (2026-10-04, Subm. \#6783).
\item
  \textbf{La idea simple:} Ningún detective es perfecto en todo. Elegir
  un explicador es como elegir un vehículo: un auto deportivo es veloz
  pero caro; una bicicleta es barata pero lenta.
\item
  \textbf{Lo que descubrimos:} Construimos la ``Frontera de Selección''.
  Le enseñamos a las organizaciones y auditores cómo elegir la mejor
  herramienta según su necesidad:

  \begin{itemize}
  \tightlist
  \item
    Si estás en un \textbf{hospital} o un \textbf{juicio}, usa
    \textbf{SHAP} (máxima estabilidad y fidelidad).
  \item
    Si estás en una \textbf{app móvil} que necesita responder en
    milisegundos, usa \textbf{LIME} (rapidez y brevedad).
  \end{itemize}
\end{itemize}

\begin{center}\rule{0.5\linewidth}{0.5pt}\end{center}

\subsubsection{6. Papel F: ¿Los Humanos y los Robots Entienden lo
Mismo?}\label{papel-f-los-humanos-y-los-robots-entienden-lo-mismo}

\begin{itemize}
\tightlist
\item
  \textbf{Estado:} En desarrollo inicial (\emph{Journal of Computer
  Sciences Institute}).
\item
  \textbf{La idea simple:} Al final del día, las explicaciones son para
  personas de carne y hueso. ¿Realmente los humanos entienden lo mismo
  que las métricas matemáticas?
\item
  \textbf{Lo que exploramos:} Diseñamos un plan para comparar si los
  juicios de las personas (y de modelos de lenguaje usados como jueces)
  coinciden con las mediciones computacionales del benchmark.
\end{itemize}

\begin{center}\rule{0.5\linewidth}{0.5pt}\end{center}

\subsection{Acto 4: El Gran Final --- La Narrativa para la Defensa
Doctoral}\label{acto-4-el-gran-final-la-narrativa-para-la-defensa-doctoral}

Esta historia une todo el trabajo de investigación en un mensaje claro y
motivador para la presentación de la tesis:

\begin{enumerate}
\def\labelenumi{\arabic{enumi}.}
\tightlist
\item
  \textbf{El Problema:} La Inteligencia Artificial es poderosa pero
  opaca. Necesitamos explicaciones.
\item
  \textbf{La Complicación:} Las herramientas de explicación existen,
  pero sin evaluación rigurosa pueden dar falsas ilusiones de
  transparencia.
\item
  \textbf{La Solución (FOM-7):} Creamos un protocolo estandarizado de 7
  puertas para auditar las explicaciones de forma reproducible.
\item
  \textbf{La Evidencia (Papeles A-E):} Probamos, comparamos y medimos
  los límites reales de herramientas como LIME, SHAP, Anchors y DiCE.
\item
  \textbf{El Impacto (Papel F y Capítulos):} Le damos a la sociedad y a
  los científicos un mapa claro para pasar de la \textbf{confianza ciega
  en los algoritmos} a una \textbf{supervisión transparente, responsable
  y centrada en las personas}.
\end{enumerate}

\begin{center}\rule{0.5\linewidth}{0.5pt}\end{center}

\emph{Nota: Este documento sirve como mapa narrativo conceptual para la
presentación doctoral y para enriquecer la divulgación del trabajo de
investigación, mantenido de manera independiente al capítulo de libro.}

\end{document}
